\section{Proofs of the fluctuation and disorder-chaos corollaries}
\label{sec:consequences}

The two consequences in the introduction now follow from the main
variance asymptotic and existing comparison inequalities. We keep the
finite-size normalizations explicit.

\subsection{The SK--Curie--Weiss comparison}

\begin{proof}[Proof of Corollary~\ref{cor:skcw}]
Abbreviate $F_n^{\mathrm{CW}}=F_n^{\mathrm{CW}}(\beta,\gamma)$ and
couple it with $F_n(\beta)$ using the same disorder. For
$\Delta_n=F_n^{\mathrm{CW}}-F_n(\beta)$, \cite[Proposition~1.5]{DeyKang2026} gives
$\Delta_n\ge0$ and, for every $u\ge0$ and $n$,
\[
 \mathbb P(\Delta_n\ge u)\le C_\gamma e^{-u},
 \qquad C_\gamma=(1-\gamma)^{-1/2}.
\]
Integrating the tail yields the uniform bounds
\[
 \begin{gathered}
 \mathbb E\Delta_n^2
 =2\int_0^\infty u\,\mathbb P(\Delta_n>u)\,du
 \le2C_\gamma,\\
 \|\Delta_n-\mathbb E\Delta_n\|_2
 \le K_\gamma:=\sqrt{2C_\gamma}.
 \end{gathered}
\]
In particular the perturbed mean and variance are finite. Let
$t_n=\sqrt{\Var F_n^{\mathrm{CW}}}$. The reverse triangle inequality
in $L^2$ gives
\[
 |t_n-s_n(\beta)|
 \le\|\Delta_n-\mathbb E\Delta_n\|_2\le K_\gamma.
\]
Since $s_n(\beta)\sim\sqrt{a_2(\beta)}n^{1/6}\to\infty$, we have
$t_n/s_n(\beta)\to1$, which proves the variance asymptotic and makes
$t_n$ positive for all sufficiently large $n$. For those $n$, using
$\|X_n(\beta)\|_2=1$,
\[
 \left\|\frac{F_n^{\mathrm{CW}}-\mathbb E F_n^{\mathrm{CW}}}{t_n}
             -X_n(\beta)\right\|_2
 \le\frac{K_\gamma}{t_n}
   +\left|\frac{s_n(\beta)}{t_n}-1\right|
 \le\frac{2K_\gamma}{t_n}\longrightarrow0.
\]
Theorem~\ref{thm:main} and Slutsky's theorem now give the asserted law.
\end{proof}

\subsection{Disorder chaos in the zero-field model}

\begin{proof}[Proof of Corollary~\ref{cor:disorder-chaos}]
For the Gaussian field in \eqref{eq:model},
\[
 \rho(\sigma,\tau):=\Cov(H_n(\sigma),H_n(\tau))
 =\frac1n\sum_{i<j}(\sigma_i\tau_i)(\sigma_j\tau_j)
 =\frac{nR^2-1}{2}.
\]
Writing a single subscript for one Gibbs average, set
\[
 q(t)=\mathbb E\langle\rho\rangle_{g,g^t}
 =\frac1n\sum_{i<j}
 \mathbb E\big[\langle\sigma_i\sigma_j\rangle_g
               \langle\sigma_i\sigma_j\rangle_{g^t}\big].
\]
These Gibbs means are smooth with all derivatives bounded. Applying
\cite[Lemma~3.3]{Chatterjee2009} to them at its two-sided time $t/2$
shows that $q$ is nonnegative and nonincreasing. The variance identity
\cite[Theorem~3.8]{Chatterjee2009}, applied to $F_n/\beta$, gives
\[
 \frac{\Var F_n(\beta)}{\beta^2}
 =\int_0^\infty e^{-s}q(s)\,ds
 \ge (1-e^{-t})q(t).
\]
Since $q(t)=(n\mathbb E\langle R^2\rangle_{g,g^t}-1)/2$, this proves
\eqref{eq:chaos-ou}.

For replacement, apply \cite[Theorem~3.14]{Chatterjee2009} directly to
the $M$ independent coordinates, regarding
$f=F_n(\beta)/(\beta\sqrt n)$ as a function of $g$. For
$\partial_{ij}=\partial/\partial g_{ij}$, throughout $\mathbb R^M$,
\[
 \partial_{ij}f=n^{-1}\langle\sigma_i\sigma_j\rangle_g,\qquad
 |\partial_{ij}f|\le n^{-1},\qquad
 |\partial_{ij}^2f|\le\beta n^{-3/2}.
\]
The Gaussian third absolute difference moment is bounded, so the
theorem's remainder is $O(M\beta n^{-5/2})=O_\beta(n^{-1/2})$,
uniformly in $k$. Its left side and variance term give
\[
 \begin{aligned}
 \frac12\left(\mathbb E\langle R^2\rangle_{g,g^A}-\frac1n\right)
 &=\mathbb E\sum_{i<j}\partial_{ij}f(g)\partial_{ij}f(g^A)\\
 &\le\frac{M+1}{k+1}\frac{\Var F_n(\beta)}{\beta^2n}
       +O_\beta(n^{-1/2}).
 \end{aligned}
\]
This proves \eqref{eq:chaos-replacement}. Finally, the variance
asymptotic in the introduction gives $\Var F_n(\beta)=O_\beta(n^{1/3})$.
Use $1-e^{-t}\ge(1-e^{-1})\min\{t,1\}$ for the first scale and
$(M+1)/(k+1)\le M/k$ for $k\ge1$ for the second. The stated sequence
conditions make the resulting bounds tend to zero.
\end{proof}
